\documentclass[12pt]{article}
\usepackage{fullpage}
\usepackage{amsmath,amssymb,amsthm}

\newtheorem{lemma}{Lemma}
\newtheorem{proposition}{Proposition}
\newcommand{\R}{\mathbb R}
\newcommand{\E}{\mathbb E}
\newcommand{\1}{\mathbf 1}
\newcommand{\Var}{\operatorname{Var}}
\newcommand{\TV}{\operatorname{TV}}
\newcommand{\Law}{\operatorname{Law}}
\newcommand{\calF}{\mathcal F}
\newcommand{\Ent}{\operatorname{Ent}}

\begin{document}

\begin{center}
{\bf Uniqueness of the invariant measure}
\end{center}

\section*{Problem statement and interpretation}

Put \(E=C_0([0,1])\), \(H=L^2(0,1)\), and let \(S_t\) be the Dirichlet heat
semigroup generated by \(\frac12\Delta\).  I use the Dirichlet version of the
Athreya--Butkovsky--L\^e--Mytnik construction in precisely the following
form.  For \(b_n=G_{1/n}\), the regularized equation has the mild solution
\[
  u_t^n=S_t x+V_t+K_t^n,\qquad
  K_t^n=\int_0^tS_{t-r}b_n(u_r^n)\,dr,
\]
where \(V_t=\int_0^tS_{t-r}\,dW_r\), and, for every \(t>0\) and
\(\Phi\in C_b(E)\),
\[
        P_t^n\Phi(x):=\E\Phi(u_t^n)\longrightarrow P_t\Phi(x).
        \tag{1}
\]
This is the weak convergence assumption in the question, written on the
Polish state space \(E\).

\begin{lemma}[A uniform regularized drift estimate]\label{lem:Hestimate}
For each \(T<\infty\) there is \(C_T<\infty\), independent of \(n\) and of
\(x\in E\), such that
\[
 \|K_t^n-S_{t-s}K_s^n\|_{L^2(\Omega;H)}
      \le C_T |t-s|^{3/4},\qquad 0\le s<t\le T .             \tag{2}
\]
\end{lemma}

\begin{proof}
Let \(D^n_{s,t}=K_t^n-S_{t-s}K_s^n\).  Write
\[
        \rho_\tau(y)=\int_0^\tau p_{2a}(y,y)\,da,\qquad
        d(y)=y\wedge(1-y).
\]
The standard one-dimensional Dirichlet kernel estimate
\(p_a(y,y)\asymp a^{-1/2}(1\wedge d(y)^2/a)\) gives
\[
        \rho_\tau(y)\asymp \sqrt{\tau}\wedge d(y).             \tag{3}
\]
Hence
\[
 \int_0^1\rho_\tau(y)^{-1/2}dy\le C\tau^{-1/4},\qquad
 \sup_z\int_0^1p_a(z,y)\rho_\tau(y)^{-1/2}dy
     \le C(a^{-1/4}+\tau^{-1/4}).                            \tag{4}
\]
Indeed, (3) reduces the first estimate to integration of
\((\sqrt{\tau}\wedge d)^{-1/2}\).  For the second, use
\(\rho_\tau^{-1/2}\le C(\tau^{-1/4}+d^{-1/2})\) and the boundary bound
\[
p_a(z,y)\le C a^{-1/2}\bigl(1\wedge d(z)d(y)/a\bigr)
              e^{-|z-y|^2/(Ca)},
\]
whose convolution with \(d(y)^{-1/2}\) is \(O(a^{-1/4})\) by the change of
variables \(y=z+\sqrt a\,r\) near the boundary.  Also
\(\rho_\tau^{-1}\le C(\tau^{-1/2}+d^{-1})\); Hardy's inequality and
\(\|\partial_xS_\tau f\|_2\le C\tau^{-1/2}\|f\|_2\) give
\[
        \|\rho_\tau^{-1}S_\tau f\|_2\le C\tau^{-1/2}\|f\|_2 .
        \tag{5}
\]

For \(s<t\), freeze the dynamics at time \(s\):
\[
 A^n_{s,t}=\int_s^tS_{t-r}
 b_n\!\left(S_{r-s}u_s^n+V_r-S_{r-s}V_s\right)\,dr .          \tag{6}
\]
For fixed \(n\), \(D^n\) is the semigroup-sewn limit of \(A^n\).  More
explicitly, if \(\Pi\) is a partition of \([s,t]\), then
\[
 \sum_{[a,b]\in\Pi}S_{t-b}A^n_{a,b}\to D^n_{s,t}\quad\hbox{in }L^2(H).
 \tag{7}
\]
To see this, compare \(A^n_{a,b}\) with
\(\int_a^bS_{b-r}b_n(u_r^n)\,dr\).  Since \(b_n\) is Lipschitz and bounded,
\[
 \Big\|A^n_{a,b}-\int_a^bS_{b-r}b_n(u_r^n)\,dr\Big\|_{L^2H}
 \le \operatorname{Lip}(b_n)\|b_n\|_\infty (b-a)^2,
\]
and summing gives an error bounded by a constant times
\((t-s)|\Pi|\).

We next estimate \(A^n\), uniformly in \(n\).  If
\(Z_r=S_{r-s}u_s^n+V_r-S_{r-s}V_s\), then, conditionally on \(\calF_s\),
\(Z_r(y)\) is Gaussian with variance \(\rho_{r-s}(y)\).  Since \(b_n\) is a
probability density, convolution with a Gaussian of variance \(\sigma^2\) has
supremum at most \((2\pi\sigma^2)^{-1/2}\).  Conditioning first at time \(r\)
therefore gives, for \(s<r<q\),
\[
 \E_s\,b_n(Z_r(y))b_n(Z_q(z))
 \le C\rho_{r-s}(y)^{-1/2}\rho_{q-r}(z)^{-1/2}.              \tag{8}
\]
Using the semigroup identity
\(\int p_{t-r}(x,y)p_{t-q}(x,z)dx=p_{2t-r-q}(y,z)\) and (4),
\[
\begin{aligned}
 \E_s\|A^n_{s,t}\|_2^2
 &\le C\!\int_s^t\!\!\int_r^t
      \bigl((2t-r-q)^{-1/4}+(r-s)^{-1/4}\bigr)(q-r)^{-1/4}\,dq\,dr  \\
 &\le C|t-s|^{3/2}.
\end{aligned}                                                \tag{9}
\]
Thus \(\|A^n_{s,t}\|_{L^2H}\le C|t-s|^{3/4}\).

Let
\(\delta A^n_{s,u,t}=A^n_{s,t}-S_{t-u}A^n_{s,u}-A^n_{u,t}\).  For
\(r>u\) the two frozen arguments differ by \(h=S_{r-u}D^n_{s,u}\).  Conditional
on \(\calF_u\), the remaining noise has variance \(\rho_{r-u}(y)\).  Moreover
\((b_n*\varphi_\sigma)'\) is bounded by \(C\sigma^{-2}\) uniformly in \(n\);
hence
\[
 \left|\E_u\{b_n(Z_r(y))-b_n(Z_r(y)+h(y))\}\right|
      \le C\rho_{r-u}(y)^{-1}|h(y)|.
\]
By (5),
\[
        \|\E_u\delta A^n_{s,u,t}\|_2
        \le C(t-u)^{1/2}\|D^n_{s,u}\|_2 .                    \tag{10}
\]

We use the following elementary Hilbert-valued sewing fact.  Let \(S_t\) be a
contraction semigroup and \(A_{s,t}\) an adapted two-parameter family.  If, for
dyadic midpoints \(u\),
\[
\|\E_s\delta A_{s,u,t}\|_{L^2H}\le \Gamma_1|t-s|^{1+\varepsilon_1},\quad
\|\delta A_{s,u,t}-\E_s\delta A_{s,u,t}\|_{L^2H}
        \le \Gamma_2|t-s|^{1/2+\varepsilon_2},
\]
then the dyadic sums \(\sum S_{t-t_{i+1}}A_{t_i,t_{i+1}}\) converge in
\(L^2(H)\) and their limit \({\mathcal A}_{s,t}\) satisfies
\[
 \|{\mathcal A}_{s,t}-A_{s,t}\|_{L^2H}
 \le C(\Gamma_1|t-s|^{1+\varepsilon_1}
       +\Gamma_2|t-s|^{1/2+\varepsilon_2}).                  \tag{11}
\]
Indeed, the difference between two consecutive dyadic sums is a sum of
\(S_{t-b}\delta A_{a,u,b}\).  The conditional-mean parts form a summable
geometric series.  The centered parts are orthogonal Hilbert-valued martingale
differences, so their \(L^2\)-norm is controlled by the square-sum
\(2^m(2^{-m}|t-s|)^{1+2\varepsilon_2}\).  This is exactly the
\(L^2(H)\) case of the stochastic sewing lemma.

For fixed \(n\), define
\[
Y_n(h)=\sup_{0<t-s\le h}
       \frac{\|D^n_{s,t}\|_{L^2H}}{|t-s|^{3/4}}<\infty .
\]
Using (10) and the tower property,
\[
 \|\E_a\delta A^n_{a,b,c}\|_{L^2H}
 \le C\,Y_n(h)(b-a)^{3/4}(c-b)^{1/2}
 \le C\,Y_n(h)(c-a)^{5/4}
\]
on intervals of length at most \(h\).  Also (9) gives
\(\|\delta A^n_{a,b,c}-\E_a\delta A^n_{a,b,c}\|_{L^2H}\le C(c-a)^{3/4}\).
Combining (7), (9), and (11),
\[
 \|D^n_{s,t}\|_{L^2H}
 \le C|t-s|^{3/4}+C\,Y_n(h)|t-s|^{5/4}.
\]
Thus \(Y_n(h)\le C+Ch^{1/2}Y_n(h)\).  Choose \(h\) so that the last term is
absorbed.  Longer intervals are split into pieces of length at most \(h\);
using \(D^n_{s,t}=S_{t-r}D^n_{s,r}+D^n_{r,t}\) yields (2), with a constant
depending only on \(T\).
\end{proof}

\section*{Reversibility and a spectral gap}

Let \(\gamma\) be Brownian bridge measure on \(E\), i.e. the centered Gaussian
measure with covariance \((-\Delta_D)^{-1}\).  Define
\[
  \Lambda(v)=\int_0^1\1_{\{v(r)>0\}}\,dr,\qquad
  \pi(dv)=Z^{-1}e^{2\Lambda(v)}\,\gamma(dv).                 \tag{12}
\]
For \(B_n'=b_n\), \(B_n(a)=\int_{-\infty}^ab_n(r)dr\), put
\[
   \pi_n(dv)=Z_n^{-1}\exp\left(2\int_0^1B_n(v(r))\,dr\right)\gamma(dv).
\]
Since \(0\le B_n\le1\), \(B_n(a)\to\1_{\{a>0\}}\) for \(a\ne0\), and a
Brownian bridge spends zero Lebesgue time at a fixed level,
\[
        \|\pi_n-\pi\|_{\TV}\to0 .                            \tag{13}
\]
All the densities \(d\pi_n/d\gamma\) and \(d\pi/d\gamma\) are bounded above
and below by deterministic positive constants.

For fixed \(n\), \(v\mapsto b_n(v)\) is a bounded globally Lipschitz map
\(H\to H\).  The regularized SPDE is the gradient diffusion for
\[
        {\cal E}_n(F,G)=\frac12\int_E\langle DF,DG\rangle_H\,d\pi_n .
\]
Here this standard identification can be obtained without formal domain
arguments as follows.  Approximate the potential by the cylinder potentials
\(2\int B_n(\Pi_m v)\), where \(\Pi_m\) is the first \(m\) Dirichlet modes.
The corresponding equation has drift \(\Pi_m b_n(\Pi_m v)\); it is the product
of a finite-dimensional reversible gradient diffusion and the independent
linear Ornstein--Uhlenbeck tail, hence is reversible with respect to
\(Z_{n,m}^{-1}e^{2\int B_n(\Pi_m v)}\gamma(dv)\).  Because
\(\Pi_m v\to v\) in \(L^2\) and \(\gamma\)-a.s. uniformly, these measures
converge in total variation to \(\pi_n\).  The Lipschitz property of \(b_n\)
gives convergence of the approximating mild solutions to \(u^n\) in
\(L^2(\Omega;C([0,T];H))\), and at each positive time in \(E\).  Passing to the
limit gives, for bounded Borel \(f,g\),
\[
       \int P_t^nf\,d\pi_n=\int f\,d\pi_n,\qquad
       \int gP_t^nf\,d\pi_n=\int fP_t^ng\,d\pi_n .           \tag{14}
\]

\begin{proposition}\label{prop:gap}
The measure \(\pi\) is \(P_t\)-invariant, \(P_t\) is symmetric on
\(L^2(\pi)\), and there exists \(c>0\) such that
\[
       \Var_\pi(P_t f)\le e^{-2ct}\Var_\pi(f),
       \qquad f\in L^2(\pi),\ t\ge0 .                       \tag{15}
\]
\end{proposition}

\begin{proof}
For \(f,g\in C_b(E)\), (1), (13), (14), and bounded convergence with respect
to the uniformly bounded densities \(d\pi_n/d\gamma\) imply
\[
 \int P_tf\,d\pi=\lim_n\int P_t^nf\,d\pi_n=\int f\,d\pi
\]
and
\[
 \int gP_tf\,d\pi=\lim_n\int gP_t^nf\,d\pi_n
       =\lim_n\int fP_t^ng\,d\pi_n=\int fP_tg\,d\pi .
\]
Thus invariance and symmetry hold, first on \(C_b(E)\) and then on
\(L^2(\pi)\) by density.

The Gaussian Poincare inequality for \(\gamma\) and the uniform two-sided
bounds on \(d\pi_n/d\gamma\) give, by Holley--Stroock, a constant \(C\)
independent of \(n\) such that
\[
      \Var_{\pi_n}(F)\le C\,{\cal E}_n(F,F).
\]
Since \(P_t^n\) is the symmetric semigroup generated by \({\cal E}_n\),
\[
      \Var_{\pi_n}(P_t^n f)\le e^{-2t/C}\Var_{\pi_n}(f).
\]
Letting \(n\to\infty\) for bounded continuous \(f\), using (1) and (13), and
then extending by \(L^2(\pi)\)-density proves (15).
\end{proof}

\section*{Fixed-time absolute continuity}

\begin{lemma}\label{lem:control}
For every deterministic \(x\in E\) and \(t>0\),
\[
        P_t(x,\cdot)\ll \gamma ,
\]
and hence \(P_t(x,\cdot)\ll\pi\).
\end{lemma}

\begin{proof}
Let \(\nu_t^x=\Law(S_tx+V_t)\).  We prove first that the regularized terminal
laws have entropy, relative to \(\nu_t^x\), bounded uniformly in \(n\).  Set
\(r_k=t(1-2^{-k})\), \(\Delta_k=r_{k+1}-r_k\), and
\(D_k^n=K^n_{r_{k+1}}-S_{\Delta_k}K^n_{r_k}\).  Define the predictable
control
\[
 F_s^n=\sum_{k\ge0}\frac{\1_{(r_{k+1},r_{k+2}]}(s)}{\Delta_{k+1}}\,
        S_{s-r_{k+1}}D_k^n .                                \tag{16}
\]
By Lemma \ref{lem:Hestimate},
\[
 \E\int_0^t\|F_s^n\|_2^2ds
 \le C_t\sum_{k\ge0}\frac{\Delta_k^{3/2}}{\Delta_{k+1}}<\infty, \tag{17}
\]
uniformly in \(n\).  Its terminal effect telescopes:
\[
 \int_0^tS_{t-s}F_s^n\,ds
 =\sum_{k\ge0}S_{t-r_{k+1}}D_k^n=K_t^n                       \tag{18}
\]
in \(H\), since the partial sums equal \(S_{t-r_{N+1}}K^n_{r_{N+1}}\) and
\(K^n\) is continuous.  If \(F\in L^2(0,t;H)\), then
\(\int_0^tS_{t-s}F_sds\in H^\alpha_0(0,1)\subset E\) for every
\(\alpha\in(1/2,1)\); this follows in the Dirichlet basis from
\[
 \sum_k\lambda_k^\alpha\Big|\int_0^te^{-\lambda_k(t-s)/2}F_k(s)ds\Big|^2
 \le \sum_k\lambda_k^{\alpha-1}\int_0^t|F_k(s)|^2ds
 \le \lambda_1^{\alpha-1}\int_0^t\|F_s\|_2^2ds .
\]
Thus all terminal variables below are \(E\)-valued.

Let \(\tau_N=\inf\{s:\int_0^s\|F_r^n\|_2^2dr>N\}\) and
\(F^{n,N}=F^n\1_{[0,\tau_N]}\).  For
\[
 X_N=S_tx+V_t+\int_0^tS_{t-s}F_s^{n,N}ds ,
\]
Novikov's condition holds.  Hilbert-space Girsanov says that the law on the
noise path space of \(W+\int_0^\cdot F_s^{n,N}ds\) has entropy at most
\(\frac12\E\int_0^t\|F_s^{n,N}\|_2^2ds\) relative to the law of \(W\).
Applying the terminal map \(W\mapsto S_tx+\int_0^tS_{t-s}dW_s\) and the
data-processing inequality,
\[
 \Ent(\Law(X_N)\mid\nu_t^x)
 \le \frac12\,\E\int_0^t\|F_s^{n,N}\|_2^2ds
 \le \frac12\,\E\int_0^t\|F_s^n\|_2^2ds .                   \tag{19}
\]
On \(\{\tau_N>t\}\), (18) gives \(X_N=u_t^n\), while
\(\mathbb P(\tau_N\le t)\le N^{-1}\E\int_0^t\|F_s^n\|_2^2ds\to0\).  Hence
\(X_N\to u_t^n\) in probability in \(E\).  Lower semicontinuity of relative
entropy on the Polish space \(E\), together with (17)--(19), yields
\[
        \sup_n\Ent(P_t^n(x,\cdot)\mid\nu_t^x)<\infty .       \tag{20}
\]
Using (1) and lower semicontinuity once more,
\(\Ent(P_t(x,\cdot)\mid\nu_t^x)<\infty\), so \(P_t(x,\cdot)\ll\nu_t^x\).

It remains to compare \(\nu_t^x\) with \(\gamma\).  In the eigenbasis
\(e_k(r)=\sqrt2\sin(k\pi r)\), the covariance of \(\nu_t^x\) has eigenvalues
\[
        q_k(t)=\frac{1-e^{-k^2\pi^2t}}{k^2\pi^2},
\]
whereas \(\gamma\) has eigenvalues \((k^2\pi^2)^{-1}\).  The ratios
\(1-e^{-k^2\pi^2t}\) are bounded away from zero and differ from \(1\) by a
square-summable sequence.  Also \(S_tx\in H^1_0(0,1)\), the Cameron--Martin
space of \(\gamma\).  The Feldman--Hajek theorem for Gaussian measures on the
Banach space \(E\) therefore gives \(\nu_t^x\sim\gamma\).  Thus
\(P_t(x,\cdot)\ll\gamma\sim\pi\).
\end{proof}

\section*{Conclusion}

Let \(\eta\) be an invariant probability measure.  If \(\pi(A)=0\), then
Lemma \ref{lem:control} gives \(P_t(y,A)=0\) for every \(y\in E\), and hence
\[
        \eta(A)=\int_E P_t(y,A)\,\eta(dy)=0.
\]
Thus \(\eta\ll\pi\); write \(d\eta=h\,d\pi\).

The detailed-balance identity for \(P_t\) extends from bounded continuous to
bounded Borel functions by a monotone-class argument applied to the finite
measure \(\pi(dx)P_t(x,dy)\) on \(E\times E\).  Hence, for every bounded
Borel \(f\),
\[
 \int f h\,d\pi=\int P_tf\,d\eta
      =\int P_tf\,h\,d\pi=\int f\,P_t h\,d\pi ,
\]
so \(P_t h=h\) in \(L^1(\pi)\).  For \(a>0\), let \(h_a=h\wedge a\).  Since
\(r\mapsto r\wedge a\) is concave,
\[
        P_t h_a\le (P_t h)\wedge a=h_a .
\]
The two sides have the same \(\pi\)-integral; therefore \(P_t h_a=h_a\).
Now \(h_a\in L^2(\pi)\), and the spectral gap (15) implies
\(\Var_\pi(h_a)=0\).  Thus \(h_a\) is constant for every \(a\).  Letting
\(a\uparrow\infty\) and using \(\int h\,d\pi=1\), we get \(h=1\) \(\pi\)-a.s.

Therefore \(P_t\) has at most one invariant probability measure on \(E\).  In
fact the invariant measure is the Gibbs measure \(\pi\) in (12).

\begin{thebibliography}{9}

\bibitem{ABLM}
S. Athreya, O. Butkovsky, K. L\^e and L. Mytnik,
\emph{Well-posedness of stochastic heat equation with distributional drift
and skew stochastic heat equation}, Comm. Pure Appl. Math. 77 (2024),
2708--2777.

\bibitem{LeBanach}
K. L\^e,
\emph{Stochastic sewing in Banach spaces}, Electron. J. Probab. 28 (2023),
paper no. 26, 1--22.

\bibitem{LeSSL}
K. L\^e,
\emph{A stochastic sewing lemma and applications}, Electron. J. Probab. 25
(2020), paper no. 38, 1--55.

\bibitem{Bogachev}
V. I. Bogachev,
\emph{Gaussian Measures}, American Mathematical Society, 1998.

\bibitem{DPZ}
G. Da Prato and J. Zabczyk,
\emph{Ergodicity for Infinite Dimensional Systems}, Cambridge University
Press, 1996.

\bibitem{MR}
Z.-M. Ma and M. R\"ockner,
\emph{Introduction to the Theory of (Non-Symmetric) Dirichlet Forms},
Springer, 1992.

\bibitem{BZ}
S. K. Bounebache and L. Zambotti,
\emph{A skew stochastic heat equation}, J. Theoret. Probab. 27 (2014),
168--201.

\end{thebibliography}

\end{document}
